% Code inventory for the Team 54 WCA-inspired application (src/ plus scripts/).
% Lists every project .py file, its path, what it contains, and its exact
% function signatures. Virtualenv and cached files are excluded on purpose.
%
% Build: pdflatex code_inventory.tex  (run twice for the table of contents)

\documentclass[11pt,a4paper]{article}

\usepackage[T1]{fontenc}
\usepackage[utf8]{inputenc}
\usepackage{underscore}
\usepackage[margin=2.2cm]{geometry}
\usepackage{longtable}
\usepackage{booktabs}
\usepackage{hyperref}
\usepackage{xcolor}
\usepackage{listings}

\lstset{
  language=Python,
  basicstyle=\ttfamily\small,
  columns=fullflexible,
  keepspaces=true,
  breaklines=true,
  breakatwhitespace=false,
  frame=single,
  framesep=4pt,
  rulecolor=\color{black!30},
  backgroundcolor=\color{black!3},
  aboveskip=4pt,
  belowskip=8pt,
  showstringspaces=false
}

\newcommand{\filepath}[1]{\texttt{\detokenize{#1}}}
\newcommand{\role}[1]{\par\smallskip\noindent\textbf{Role.} #1\par}
\newcommand{\filesize}[1]{\par\smallskip\noindent\textbf{Size.} #1 lines\par}
\newcommand{\contents}[1]{\par\smallskip\noindent\textbf{Contains.} #1\par\smallskip}

\title{Code Inventory\\\large Team 54 --- WCA-inspired mini-world application}
\author{Generated from the repository at commit \texttt{818b78a}}
\date{18 November 2026}

\begin{document}

\maketitle

\noindent
This document inventories every Python source file that belongs to the
project (everything under \filepath{scripts/} and \filepath{src/} except the
virtualenv and \filepath{__pycache__}). For each file it gives the path, a
one-paragraph description of what the file contains, and the exact
signatures of every function defined in it. Signatures are copied verbatim
from the source, including type hints, default values and FastAPI
\texttt{Query}/\texttt{Body} declarations. Line numbers in comments
(\texttt{\# Lnn}) point at the \texttt{def} statement.

\tableofcontents
\newpage

% ---------------------------------------------------------------------------
\section*{Summary}
\addcontentsline{toc}{section}{Summary}

\begin{longtable}{@{}p{0.46\textwidth} r p{0.40\textwidth}@{}}
\toprule
\textbf{File} & \textbf{Lines} & \textbf{Kind} \\
\midrule
\endhead
\filepath{scripts/make_mock_data.py}          & 173 & Dev-only fixture generator \\
\filepath{src/app.py}                         & 19  & FastAPI application object \\
\filepath{src/config.py}                      & 7   & Environment settings \\
\filepath{src/data/__init__.py}               & 8   & Mock/real implementation switch \\
\filepath{src/data/mock.py}                   & 1159& Dev-only in-memory data layer \\
\filepath{src/data/real.py}                   & 164 & Real database layer (stubs) \\
\filepath{src/errors.py}                      & 51  & Domain errors and handlers \\
\filepath{src/routers/__init__.py}            & 0   & Empty package marker \\
\filepath{src/routers/competitions.py}        & 123 & Competition endpoints \\
\filepath{src/routers/lookups.py}             & 36  & Country and event endpoints \\
\filepath{src/routers/persons.py}             & 46  & Person endpoints \\
\filepath{src/routers/rankings.py}            & 27  & Rankings endpoint \\
\filepath{src/routers/search.py}              & 14  & Global search endpoint \\
\filepath{src/routers/stats.py}               & 29  & Statistics endpoints \\
\filepath{src/schemas.py}                     & 296 & Pydantic request/response models \\
\filepath{src/serialization.py}               & 18  & JSON-safe value converter \\
\filepath{src/tests/test_api.py}              & 229 & HTTP/API tests \\
\filepath{src/tests/test_contract.py}         & 366 & Data-layer contract tests \\
\bottomrule
\end{longtable}

\noindent
Total: 18 Python files, 2765 lines.

\newpage

% ---------------------------------------------------------------------------
\section{\filepath{scripts/make_mock_data.py}}

\noindent\textbf{Path.} \filepath{scripts/make_mock_data.py}
\filesize{173}
\role{Development-only script (deleted before submission). It reads the
delivered WCA-shaped CSVs from \filepath{data/} and writes one JSON fixture per
table into \filepath{src/mock\_data/}. A module-level \texttt{SPEC} dictionary
declares, for each table, the exact column set and a coercion type
(\texttt{INT}, \texttt{FLOAT}, \texttt{BOOL}, \texttt{STR}). Confidential
columns are not copied, so they can never reach the UI. Run with
\texttt{uv run -{}-no-project python scripts/make\_mock\_data.py}.}
\contents{The \texttt{SPEC} table/column schema, three helper functions and a
\texttt{\_\_main\_\_} guard.}

\begin{lstlisting}
def coerce(raw, kind):                                    # L140

def build(table, columns):                                # L152

def main():                                               # L161
\end{lstlisting}

% ---------------------------------------------------------------------------
\section{\filepath{src/app.py}}

\noindent\textbf{Path.} \filepath{src/app.py}
\filesize{19}
\role{FastAPI entry point. Creates the \texttt{FastAPI} application, registers
the error handlers from \filepath{errors.py}, includes the six routers
(\texttt{search}, \texttt{competitions}, \texttt{persons}, \texttt{rankings},
\texttt{lookups}, \texttt{stats}) under the \texttt{/api} prefix, and finally
mounts \filepath{static/} with \texttt{html=True} so that \texttt{index.html}
is served at \texttt{/} and \filepath{404.html} catches unknown paths. The
static mount is deliberately last so the API routes win.}
\contents{Module-level application wiring only: \texttt{BASE}, \texttt{app} and
the router-inclusion loop.}

\noindent\textbf{Defines no functions.}

% ---------------------------------------------------------------------------
\section{\filepath{src/config.py}}

\noindent\textbf{Path.} \filepath{src/config.py}
\filesize{7}
\role{Database connection settings read from environment variables with safe
defaults: \texttt{DB\_HOST} (\texttt{localhost}), \texttt{DB\_PORT}
(\texttt{3306}), \texttt{DB\_USER} (\texttt{root}), \texttt{DB\_PASSWORD}
(empty) and \texttt{DB\_NAME} (\texttt{miniworld}). Every other module takes
credentials from here, never from literals.}
\contents{Module-level constants only.}

\noindent\textbf{Defines no functions.}

% ---------------------------------------------------------------------------
\section{\filepath{src/data/__init__.py}}

\noindent\textbf{Path.} \filepath{src/data/\_\_init\_\_.py}
\filesize{8}
\role{The single switch between the mock and the real data layer. It imports
\texttt{data.mock} when \texttt{USE\_MOCK=1} and \texttt{data.real} otherwise,
then re-exports every public attribute of the chosen module through a
module-level \texttt{\_\_getattr\_\_}. Callers write \texttt{import data} and
\texttt{data.list\_competitions(...)}; the implementation is invisible to
routers and tests.}
\contents{One dispatch variable and one dunder accessor.}

\begin{lstlisting}
def __getattr__(name):                                    # L7
\end{lstlisting}

% ---------------------------------------------------------------------------
\section{\filepath{src/data/real.py}}

\noindent\textbf{Path.} \filepath{src/data/real.py}
\filesize{164}
\role{The real data layer. This is the only file allowed to contain raw SQL
(or MQL/Cypher) and the only place where driver exceptions are translated into
the domain errors \texttt{NotFound}, \texttt{Conflict} and \texttt{BadInput}.
The signatures are final and match \filepath{docs/feature\_specs.md}; the
bodies currently raise \texttt{NotImplementedError} and will be filled in once
the database teammate delivers the schema. It also declares the three sort
whitelists (\texttt{COMPETITION\_SORTS}, \texttt{PERSON\_SORTS},
\texttt{RESULT\_SORTS}) that keep \texttt{ORDER BY} clauses out of user
control.}
\contents{Three whitelist dictionaries, 16 read functions and 4 write
functions, all typed.}

\subsection*{Reads}
\begin{lstlisting}
def search_all(q: str, limit: int = 5) -> dict:                                 # L32

def list_competitions(
    q: Optional[str] = None,
    country_id: Optional[str] = None,
    continent_id: Optional[str] = None,
    year: Optional[int] = None,
    status: Optional[str] = None,
    event_id: Optional[str] = None,
    sort: str = "date_desc",
    page: int = 1,
    page_size: int = 20,
) -> dict:                                                                      # L36

def get_competition(competition_id: str) -> dict:                               # L50

def list_competition_results(
    competition_id: str,
    event_id: Optional[str] = None,
    round_number: Optional[int] = None,
    tab: str = "podiums",
    sort: Optional[str] = None,
    page: int = 1,
    page_size: int = 20,
) -> dict:                                                                      # L54

def list_persons(
    q: Optional[str] = None,
    country_id: Optional[str] = None,
    sort: str = "name",
    page: int = 1,
    page_size: int = 20,
) -> dict:                                                                      # L66

def get_person(person_ref: Union[int, str]) -> dict:                            # L76

def list_person_results(
    person_ref: Union[int, str],
    event_id: Optional[str] = None,
    page: int = 1,
    page_size: int = 20,
) -> dict:                                                                      # L85

def list_rankings(
    event_id: str,
    kind: str = "single",
    country_id: Optional[str] = None,
    continent_id: Optional[str] = None,
    page: int = 1,
    page_size: int = 100,
) -> dict:                                                                      # L94

def list_countries() -> list:                                                   # L105

def get_country(country_id: str) -> dict:                                       # L109

def list_events(include_retired: bool = False) -> list:                         # L113

def list_registrations(
    competition_id: str,
    status: Optional[str] = None,
    page: int = 1,
    page_size: int = 100,
) -> dict:                                                                      # L117

def get_registration(competition_id: str, registrant_no: int) -> dict:          # L126

def stats_medals() -> list:                                                     # L130

def stats_competitions_by_country_year(
    year_from: Optional[int] = None, year_to: Optional[int] = None
) -> list:                                                                      # L134

def stats_records() -> list:                                                    # L140
\end{lstlisting}

\subsection*{Writes}
\begin{lstlisting}
def create_registration(
    competition_id: str, person_no: int, event_ids: list, guests: int = 0
) -> dict:                                                                      # L147

def set_registration_status(
    competition_id: str, registrant_no: int, status: str
) -> dict:                                                                      # L153

def delete_registration(competition_id: str, registrant_no: int) -> None:       # L159

def update_person_country(person_ref: Union[int, str], country_id: str) -> dict: # L163
\end{lstlisting}

% ---------------------------------------------------------------------------
\section{\filepath{src/data/mock.py}}

\noindent\textbf{Path.} \filepath{src/data/mock.py}
\filesize{1159}
\role{Development-only in-memory implementation of the exact same contract as
\filepath{real.py} (deleted before submission). It loads the JSON fixtures from
\filepath{src/mock\_data/} with \texttt{lru\_cache}, mutates them in place for
writes, and may filter, sort and paginate in Python --- something the real
layer must never do. It reproduces pagination bounds, the sort whitelist, tie
handling in rankings, waiting-list promotion and every \texttt{NotFound} /
\texttt{Conflict} / \texttt{BadInput} case, so the routers and the tests run
unchanged against either layer.}
\contents{53 module-level functions plus 3 nested helpers: cached fixture
accessors, row-shaping helpers, the 20 contract functions mirroring
\filepath{real.py}, and a cache-reset hook used by the test fixtures.}

\subsection*{Fixture loading and cache control}
\begin{lstlisting}
def _load(table):                                         # L23
def _reset():                                             # L27
def _reset_derived():                                     # L1143
\end{lstlisting}

\subsection*{Generic helpers}
\begin{lstlisting}
def _as_date(value):                                      # L36
def _as_datetime(value):                                  # L42
def _competition_status(comp):                            # L48
def _matches(value, needle):                              # L60
def _page(rows, page, page_size):                          # L68
def _sort_key(sorts, sort, default):                       # L79
\end{lstlisting}

\subsection*{Cached lookup tables}
\begin{lstlisting}
def _persons():                                           # L91
def _countries():                                         # L96
def _continents():                                        # L101
def _events():                                            # L106
def _formats():                                           # L111
def _competitions():                                      # L116
def _championship_types():                                # L121
def _delegate_status():                                   # L126
def _final_rounds():                                      # L159
def _attempts_by_result():                                # L168
def _selects_by_registration():                           # L178
def _registered_counts():                                 # L188
def _person_stats():                                      # L505
\end{lstlisting}

\subsection*{Row shaping and resolution}
\begin{lstlisting}
def _resolve_person(person_ref):                          # L130
def _person_name(person_no):                              # L144
def _country_name(country_id):                            # L149
def _result_key(result):                                  # L154
def _competition_row(comp):                               # L197
def _competition_results(
    competition_id, event_id=None, round_number=None, tab="podiums"
):                                                        # L422
def _personal_records(person_no):                         # L551
def _registration_row(reg):                               # L807
def _competition_for_write(competition_id):               # L934
def _is_registration_open(comp):                          # L941
def _offered_events(competition_id):                      # L954
\end{lstlisting}

\subsection*{Contract: reads (mirrors \filepath{real.py})}
\begin{lstlisting}
def search_all(q, limit=5):                               # L219

def list_competitions(
    q=None,
    country_id=None,
    continent_id=None,
    year=None,
    status=None,
    event_id=None,
    sort="date_desc",
    page=1,
    page_size=20,
):                                                        # L262

def get_competition(competition_id):                      # L314

def list_competition_results(
    competition_id,
    event_id=None,
    round_number=None,
    tab="podiums",
    sort=None,
    page=1,
    page_size=20,
):                                                        # L462

def list_persons(q=None, country_id=None, sort="name", page=1, page_size=20):  # L520

def get_person(person_ref):                               # L631

def list_person_results(person_ref, event_id=None, page=1, page_size=20):      # L650

def list_rankings(
    event_id, kind="single", country_id=None, continent_id=None, page=1, page_size=100
):                                                        # L683

def list_countries():                                     # L770
def get_country(country_id):                              # L783
def list_events(include_retired=False):                   # L790

def list_registrations(competition_id, status=None, page=1, page_size=100):    # L825

def get_registration(competition_id, registrant_no):      # L837

def stats_medals():                                       # L847
def stats_competitions_by_country_year(year_from=None, year_to=None):          # L879
def stats_records():                                      # L902
\end{lstlisting}

\subsection*{Contract: writes}
\begin{lstlisting}
def create_registration(competition_id, person_no, event_ids, guests=0):       # L960

def set_registration_status(competition_id, registrant_no, status):            # L1022

def delete_registration(competition_id, registrant_no):                        # L1074

def update_person_country(person_ref, country_id):                             # L1117
\end{lstlisting}

\subsection*{Nested helpers}
\begin{lstlisting}
# inside get_competition
def people(pairs):                                        # L356

# inside _personal_records
def count_better(event, value, key, scope=None):          # L579

# inside list_rankings
def in_scope(row):                                        # L708
\end{lstlisting}

% ---------------------------------------------------------------------------
\section{\filepath{src/errors.py}}

\noindent\textbf{Path.} \filepath{src/errors.py}
\filesize{51}
\role{The domain error hierarchy and the three exception handlers that turn
errors into the uniform JSON body \texttt{\{"detail": ..., "code": ...\}}. The
base class \texttt{AppError} carries a \texttt{status\_code},
\texttt{detail} and machine-readable \texttt{code}; \texttt{NotFound} maps to
404, \texttt{Conflict} to 409 and \texttt{BadInput} to 422.
\texttt{RequestValidationError} is flattened into one human sentence, and an
unhandled exception is logged and returned as a generic 500 so no driver
detail escapes.}
\contents{Three exception classes plus \texttt{register\_error\_handlers},
which installs three nested handler coroutines on the app.}

\begin{lstlisting}
class AppError(Exception):                                # L10
    def __init__(self, detail: str, code: str | None = None):   # L13

class NotFound(AppError):                                 # L17
class Conflict(AppError):                                 # L21
class BadInput(AppError):                                 # L25

def register_error_handlers(app):                         # L29
    # nested inside register_error_handlers
    async def app_error(_, exc: AppError):                # L31
    async def validation_error(_, exc: RequestValidationError):  # L37
    async def fallback(_, exc: Exception):                # L46
\end{lstlisting}

% ---------------------------------------------------------------------------
\section{\filepath{src/routers/__init__.py}}

\noindent\textbf{Path.} \filepath{src/routers/\_\_init\_\_.py}
\filesize{0}
\role{Empty file that marks \filepath{routers/} as a package so
\texttt{from routers import ...} works.}
\contents{Nothing.}

\noindent\textbf{Defines no functions.}

% ---------------------------------------------------------------------------
\section{\filepath{src/routers/competitions.py}}

\noindent\textbf{Path.} \filepath{src/routers/competitions.py}
\filesize{123}
\role{All competition endpoints. Every handler validates its query parameters
at the edge with FastAPI \texttt{Query} bounds and enum hints, delegates
straight to \texttt{data.*}, and declares its response model so the OpenAPI
schema is exact. Covers the list with filters and sort, the detail page,
results with three tabs, registration listing/reading, and the three
registration writes (create with a \texttt{Location} header, status change, and
delete with a 204).}
\contents{One \texttt{APIRouter} with \texttt{tags=["competitions"]} and eight
route functions.}

\begin{lstlisting}
@router.get("/competitions", response_model=schemas.Page[schemas.CompetitionListItem])
def list_competitions(
    q: Optional[str] = Query(default=None, max_length=60),
    country_id: Optional[str] = Query(default=None, max_length=60),
    continent_id: Optional[str] = Query(default=None, max_length=20),
    year: Optional[int] = Query(default=None, ge=2003, le=2100),
    status: Optional[schemas.CompetitionStatus] = None,
    event_id: Optional[str] = Query(default=None, max_length=10),
    sort: str = Query(default="date_desc", max_length=20),
    page: int = Query(default=1, ge=1),
    page_size: int = Query(default=20, ge=1, le=100),
):                                                        # L12

@router.get("/competitions/{competition_id}", response_model=schemas.CompetitionDetail)
def get_competition(competition_id: str):                 # L37

@router.get("/competitions/{competition_id}/results",
            response_model=schemas.Page[schemas.CompetitionResultItem])
def list_competition_results(
    competition_id: str,
    tab: str = Query(default="podiums", pattern="^(podiums|all|by_person)$"),
    event_id: Optional[str] = Query(default=None, max_length=10),
    round_number: Optional[int] = Query(default=None, ge=1, le=10),
    sort: Optional[str] = Query(default=None, max_length=20),
    page: int = Query(default=1, ge=1),
    page_size: int = Query(default=20, ge=1, le=100),
):                                                        # L45

@router.get("/competitions/{competition_id}/registrations",
            response_model=schemas.Page[schemas.RegistrationItem])
def list_registrations(
    competition_id: str,
    status: Optional[schemas.RegistrationStatus] = None,
    page: int = Query(default=1, ge=1),
    page_size: int = Query(default=100, ge=1, le=100),
):                                                        # L69

@router.get("/competitions/{competition_id}/registrations/{registrant_no}",
            response_model=schemas.RegistrationItem)
def get_registration(competition_id: str, registrant_no: int):   # L83

@router.post("/competitions/{competition_id}/registrations",
             response_model=schemas.RegistrationItem,
             status_code=status.HTTP_201_CREATED)
def create_registration(
    competition_id: str, payload: schemas.RegistrationCreate, response: Response
):                                                        # L92

@router.patch("/competitions/{competition_id}/registrations/{registrant_no}",
              response_model=schemas.RegistrationChange)
def set_registration_status(
    competition_id: str, registrant_no: int, payload: schemas.RegistrationStatusUpdate
):                                                        # L109

@router.delete("/competitions/{competition_id}/registrations/{registrant_no}",
               status_code=status.HTTP_204_NO_CONTENT)
def delete_registration(competition_id: str, registrant_no: int):   # L121
\end{lstlisting}

% ---------------------------------------------------------------------------
\section{\filepath{src/routers/lookups.py}}

\noindent\textbf{Path.} \filepath{src/routers/lookups.py}
\filesize{36}
\role{Reference-data endpoints used to populate filter dropdowns and the
entity-specific country page: the country list, one country, all persons for
one country, and the event list (official only by default,
\texttt{include\_retired=true} to include historical events).}
\contents{One \texttt{APIRouter} with \texttt{tags=["lookups"]} and four route
functions.}

\begin{lstlisting}
@router.get("/countries", response_model=list[schemas.CountryItem])
def list_countries():                                     # L10

@router.get("/countries/{country_id}", response_model=schemas.CountryItem)
def get_country(country_id: str):                         # L15

@router.get("/countries/{country_id}/persons",
            response_model=schemas.Page[schemas.PersonListItem])
def list_country_persons(
    country_id: str,
    sort: str = Query(default="name", max_length=20),
    page: int = Query(default=1, ge=1),
    page_size: int = Query(default=20, ge=1, le=100),
):                                                        # L23

@router.get("/events", response_model=list[schemas.EventItem])
def list_events(include_retired: bool = False):           # L35
\end{lstlisting}

% ---------------------------------------------------------------------------
\section{\filepath{src/routers/persons.py}}

\noindent\textbf{Path.} \filepath{src/routers/persons.py}
\filesize{46}
\role{Person endpoints: the searchable, sortable person list; the person
detail page (personal records with world/continental/national ranks plus
competition and podium counts); the person's full result history; and the
secondary write that changes a person's represented country. The
\texttt{\{person\_ref\}} path segment accepts either a numeric
\texttt{Person\_No} or a WCA ID, which the data layer resolves.}
\contents{One \texttt{APIRouter} with \texttt{tags=["persons"]} and four route
functions.}

\begin{lstlisting}
@router.get("/persons", response_model=schemas.Page[schemas.PersonListItem])
def list_persons(
    q: Optional[str] = Query(default=None, max_length=60),
    country_id: Optional[str] = Query(default=None, max_length=60),
    sort: str = Query(default="name", max_length=20),
    page: int = Query(default=1, ge=1),
    page_size: int = Query(default=20, ge=1, le=100),
):                                                        # L12

@router.get("/persons/{person_ref}", response_model=schemas.PersonDetail)
def get_person(person_ref: str):                          # L25

@router.get("/persons/{person_ref}/results",
            response_model=schemas.Page[schemas.PersonResultItem])
def list_person_results(
    person_ref: str,
    event_id: Optional[str] = Query(default=None, max_length=10),
    page: int = Query(default=1, ge=1),
    page_size: int = Query(default=20, ge=1, le=100),
):                                                        # L33

@router.put("/persons/{person_ref}/country", response_model=schemas.PersonDetail)
def update_person_country(person_ref: str, payload: schemas.PersonCountryUpdate):  # L45
\end{lstlisting}

% ---------------------------------------------------------------------------
\section{\filepath{src/routers/rankings.py}}

\noindent\textbf{Path.} \filepath{src/routers/rankings.py}
\filesize{27}
\role{Single endpoint for event rankings. Ranks are computed in the data layer
over the in-scope rows (world, or narrowed by country or continent), with ties
sharing a rank. Retired events and an invalid \texttt{kind} are rejected by the
data layer with \texttt{BadInput}.}
\contents{One \texttt{APIRouter} with \texttt{tags=["rankings"]} and one route
function.}

\begin{lstlisting}
@router.get("/rankings/{event_id}", response_model=schemas.Page[schemas.RankingItem])
def list_rankings(
    event_id: str,
    kind: schemas.ResultKind = "single",
    country_id: Optional[str] = Query(default=None, max_length=60),
    continent_id: Optional[str] = Query(default=None, max_length=20),
    page: int = Query(default=1, ge=1),
    page_size: int = Query(default=100, ge=1, le=100),
):                                                        # L12
\end{lstlisting}

% ---------------------------------------------------------------------------
\section{\filepath{src/routers/search.py}}

\noindent\textbf{Path.} \filepath{src/routers/search.py}
\filesize{14}
\role{Global header search. One required query string \texttt{q} (1--60
characters) and a per-group result limit. Returns a grouped payload with
matching competitions and persons; an all-whitespace query is rejected by the
data layer.}
\contents{One \texttt{APIRouter} with \texttt{tags=["search"]} and one route
function.}

\begin{lstlisting}
@router.get("/search", response_model=schemas.SearchResults)
def search(
    q: str = Query(min_length=1, max_length=60),
    limit: int = Query(default=5, ge=1, le=20),
):                                                        # L10
\end{lstlisting}

% ---------------------------------------------------------------------------
\section{\filepath{src/routers/stats.py}}

\noindent\textbf{Path.} \filepath{src/routers/stats.py}
\filesize{29}
\role{Statistics endpoints backing the stats page: the medal table (gold,
silver and bronze counts per country from final-round podiums), competitions
per country per year with an optional year range, and record counts per event
and record code. All aggregation happens in the data layer, never in the
router.}
\contents{One \texttt{APIRouter} with \texttt{tags=["stats"]} and three route
functions.}

\begin{lstlisting}
@router.get("/stats/medals", response_model=list[schemas.MedalRow])
def stats_medals():                                       # L12

@router.get("/stats/competitions-by-country", response_model=list[schemas.CountryYearRow])
def stats_competitions_by_country(
    year_from: Optional[int] = Query(default=None, ge=2003, le=2100),
    year_to: Optional[int] = Query(default=None, ge=2003, le=2100),
):                                                        # L20

@router.get("/stats/records", response_model=list[schemas.RecordRow])
def stats_records():                                      # L28
\end{lstlisting}

% ---------------------------------------------------------------------------
\section{\filepath{src/schemas.py}}

\noindent\textbf{Path.} \filepath{src/schemas.py}
\filesize{296}
\role{Every Pydantic model used at the API edge. Response models pin the exact
JSON shape each endpoint may return (which is also the confidentiality
guarantee: confidential fields such as date of birth and email are simply not
part of any model); request models enforce types, ranges and list lengths before
the data layer runs. The generic \texttt{Page[T]} wrapper is the shared list
envelope \texttt{\{items, total, page, page\_size\}}. Five \texttt{Literal}
aliases define the allowed enum values.}
\contents{5 type aliases, a \texttt{TypeVar}, and 28 model classes. The file
defines no functions.}

\begin{lstlisting}
RegistrationStatus   = Literal["accepted", "waiting_list", "pending", "cancelled", "rejected"]
CompetitionStatus    = Literal["upcoming", "in_progress", "past", "cancelled"]
ResultKind           = Literal["single", "average"]
DelegateStatus       = Literal["senior", "regional", "full", "junior", "trainee"]
T = TypeVar("T")

class Page(BaseModel, Generic[T]):                        # L14   items, total, page, page_size
class ErrorBody(BaseModel):                               # L21   detail, code
class PersonListItem(BaseModel):                          # L29   person_no, wca_id, name, country_id,
                                                          #       country_name, competitions_count, podiums_count
class PersonalRecord(BaseModel):                          # L39   event_id, event_name, best, average,
                                                          #       world_rank, continental_rank, national_rank
class PersonDetail(BaseModel):                            # L49   person_no, wca_id, name, local_name, gender,
                                                          #       country_id, country_name, since_date,
                                                          #       delegate_status, competitions_count,
                                                          #       podiums_count, personal_records
class PersonResultItem(BaseModel):                        # L64   result_id, competition_id, competition_name,
                                                          #       event_id, round_number, position, best,
                                                          #       average, single_record, average_record,
                                                          #       country_id, solves
class EventItem(BaseModel):                               # L79   event_id, event_name, result_kind, status
class RoundItem(BaseModel):                               # L89   round_number, format_id, format_name,
                                                          #       per_attempt, cutoff_attempts, cutoff_value,
                                                          #       adv_type, adv_level, is_dual,
                                                          #       is_head_to_head, last_day
class CompetitionEvent(BaseModel):                        # L103  event_id, event_name, result_kind, rounds
class PersonRef(BaseModel):                               # L110  person_no, name, wca_id
class PodiumItem(BaseModel):                              # L116  event_id, event_name, position, person_no,
                                                          #       name, wca_id, country_id, best, average,
                                                          #       single_record, average_record
class CompetitionListItem(BaseModel):                     # L130  competition_id, name, city, country_id,
                                                          #       country_name, continent_id, start_date,
                                                          #       end_date, status, competitor_limit,
                                                          #       registered_count, main_event_id,
                                                          #       championship_type
class CompetitionDetail(CompetitionListItem):             # L146  + short_name, venue_name, street_address,
                                                          #       venue_details, latitude, longitude, website,
                                                          #       opens_at, closes_at, amount, currency_code,
                                                          #       refund_policy_percent, announced_at,
                                                          #       cancelled_at, results_posted_at, events,
                                                          #       delegates, organizers, podiums
class CompetitionResultItem(BaseModel):                   # L168  result_id, event_id, round_number, position,
                                                          #       person_no, name, wca_id, country_id, best,
                                                          #       average, single_record, average_record, solves
class RankingItem(BaseModel):                             # L187  rank, person_no, name, wca_id, country_id,
                                                          #       country_name, best, average,
                                                          #       competition_id, competition_name,
                                                          #       continental_rank, national_rank
class SearchCompetition(BaseModel):                       # L202  competition_id, name, city, country_id, start_date
class SearchPerson(BaseModel):                            # L210  person_no, wca_id, name, country_id
class SearchResults(BaseModel):                           # L217  competitions, persons
class CountryItem(BaseModel):                             # L222  country_id, iso2, name, continent_id,
                                                          #       continent_name
class MedalRow(BaseModel):                                # L230  country_id, country_name, continent_id,
                                                          #       gold, silver, bronze, total
class CountryYearRow(BaseModel):                          # L240  country_id, country_name, year,
                                                          #       competition_count
class RecordRow(BaseModel):                               # L247  event_id, event_name, kind, record_code, count
class RegistrationItem(BaseModel):                        # L258  competition_id, registrant_no, person_no, name,
                                                          #       wca_id, status, registered_at, guests,
                                                          #       guardian_consent, events
class RegistrationStatusPromotion(BaseModel):             # L271  registrant_no, person_no, name
class RegistrationChange(BaseModel):                      # L277  registration, promoted
class RegistrationCreate(BaseModel):                      # L285  person_no, events, guests
class RegistrationStatusUpdate(BaseModel):                # L291  status
class PersonCountryUpdate(BaseModel):                     # L295  country_id
\end{lstlisting}

% ---------------------------------------------------------------------------
\section{\filepath{src/serialization.py}}

\noindent\textbf{Path.} \filepath{src/serialization.py}
\filesize{18}
\role{Recursive converter used by \filepath{real.py} before returning anything:
it walks dicts, lists, tuples and sets, converts \texttt{Decimal} to
\texttt{float}, \texttt{datetime}/\texttt{date}/\texttt{time} to ISO strings,
neo4j-style objects via \texttt{to\_native()}, and Mongo \texttt{ObjectId} to
\texttt{str}, so every response is JSON-safe no matter which driver produced
it.}
\contents{One function.}

\begin{lstlisting}
def to_jsonable(o):                                       # L5
\end{lstlisting}

% ---------------------------------------------------------------------------
\section{\filepath{src/tests/test_api.py}}

\noindent\textbf{Path.} \filepath{src/tests/test\_api.py}
\filesize{229}
\role{Development-only HTTP tests run through \texttt{TestClient} with
\texttt{raise\_server\_exceptions=False} so the catch-all 500 handler can be
asserted. They check that every documented route exists in
\texttt{/openapi.json}, that reads return 200, that list responses use the
\texttt{\{items, total, page, page\_size\}} envelope, that errors are always
\texttt{\{detail, code\}}, that query validation rejects out-of-range values,
that writes return 201/200/204 (and 404 on a second delete), that blocked
deletes and country changes surface as readable 409s, that confidential fields
never leave the API, that unhandled exceptions become the generic 500 shape, and
that the static frontend is served. The autouse fixture resets the mock before
each test.}
\contents{One fixture, 17 test functions and one local helper.}

\begin{lstlisting}
@pytest.fixture(autouse=True)
def reset():                                              # L21

def test_every_documented_route_exists():                 # L29
def test_reads_return_200(url):                           # L76   (parametrized)
def test_list_envelope_shape():                           # L81
def test_error_shape_is_detail_and_code():                # L88
def test_not_found_for_each_resource():                   # L98
def test_validation_errors_are_flattened_to_one_message():# L110
def test_query_validation_bounds():                       # L118
def test_data_layer_errors_are_not_500s():                # L132
def test_create_returns_201_with_location():              # L138
def test_create_conflict_and_validation():                # L149
def test_patch_returns_registration_and_promotion():      # L168
def test_delete_returns_204_then_404():                   # L179
def test_delete_blocked_is_409_with_readable_message():   # L185
def test_update_person_country():                         # L193
def test_confidential_fields_never_leave_the_api():       # L203
def test_unhandled_error_returns_500_shape(monkeypatch):  # L211
    # nested inside test_unhandled_error_returns_500_shape
    def explode(*_args, **_kwargs):                       # L214
def test_static_frontend_is_served():                     # L226
\end{lstlisting}

% ---------------------------------------------------------------------------
\section{\filepath{src/tests/test_contract.py}}

\noindent\textbf{Path.} \filepath{src/tests/test\_contract.py}
\filesize{366}
\role{Development-only data-layer contract tests. They call \texttt{data.*}
directly and assert the shapes, pagination bounds, sort whitelisting,
confidentiality, tie handling, waiting-list promotion, and the
\texttt{NotFound} / \texttt{Conflict} / \texttt{BadInput} behaviour, so the
same file can run unchanged against the mock and later against the real layer.
The autouse fixture resets the mock before every test.}
\contents{One fixture and 39 test functions grouped by feature: search,
competitions, results, persons, rankings, lookups, registrations, stats, and
the four writes.}

\begin{lstlisting}
@pytest.fixture(autouse=True)
def reset():                                              # L18

# search
def test_search_groups_competitions_and_persons():        # L25
def test_search_by_wca_id_prefix():                       # L31
def test_search_rejects_empty_query():                    # L36
def test_search_respects_limit():                         # L41
def test_search_does_not_treat_like_wildcards_as_patterns():  # L47

# competitions
def test_list_competitions_shape_and_paging():            # L51
def test_list_competitions_filters():                     # L62
def test_list_competitions_rejects_unknown_sort():        # L79
def test_list_competitions_paging_bounds():               # L84
def test_competition_detail_has_related_data():           # L91
def test_competition_detail_never_exposes_confidential_fields():  # L100
def test_missing_competition_raises_not_found():          # L106

# competition results
def test_competition_results_tabs():                      # L115
def test_competition_result_attempts_are_ordered():       # L127
def test_competition_results_rejects_unknown_tab():       # L134

# persons
def test_persons_list_shape_and_sorting():                # L139
def test_persons_filter_by_country_and_query():           # L150
def test_person_resolves_by_number_and_wca_id():          # L158
def test_person_detail_never_exposes_confidential_fields():  # L165
def test_person_personal_records_have_ranks():            # L172
def test_missing_person_raises_not_found():               # L185
def test_person_results_for_newcomer_is_empty_not_an_error():  # L192

# rankings
def test_rankings_order_and_ties():                       # L197
def test_rankings_filter_by_country():                    # L208
def test_rankings_reject_retired_event_and_bad_kind():    # L214

# lookups
def test_countries_and_events():                          # L223

# registrations and stats
def test_registrations_list_and_filter():                 # L235
def test_stats_shapes():                                  # L245

# writes
def test_create_registration_and_selects():               # L261
def test_create_registration_duplicate_is_conflict():     # L275
def test_create_registration_validations():               # L280
def test_status_change_promotes_earliest_waiting_registration():  # L291
def test_status_change_to_waiting_list_does_not_promote_itself():  # L310
def test_status_change_without_waiting_list_promotes_nobody():     # L319
def test_status_change_validations():                     # L324
def test_delete_registration_removes_selects():           # L333
def test_delete_registration_blocked_by_results():        # L344
def test_update_person_country_sets_date():               # L350
def test_update_person_country_validations():             # L359
\end{lstlisting}

\end{document}
